% output=pdftex

\usemodule[abr-02]

\setupcolors
  [state=start]

\setuplayout
  [width=middle,height=middle,
   backspace=2cm,topspace=2cm,
   footer=0pt,
   bottomspace=2cm]

\setupwhitespace[big]

\usetypescriptfile[type-buy]

\definetypeface [mainface] [rm] [serif] [palatino] [default]
\definetypeface [mainface] [ss] [sans]  [univers]  [default]
\definetypeface [mainface] [tt] [mono]  [modern]   [default]

\setupbodyfont[mainface]

\startbuffer[colors]
\definecolor [cyan]    [c=1,m=.25,y=.25,k=.25]
\definecolor [magenta] [c=.25,m=1,y=.25,k=.25]
\definecolor [yellow]  [c=.25,m=.5,y=1,k=.25]
\definecolor [black]   [c=.25,m=.5,y=.25,k=1]
\stopbuffer

\getbuffer[colors]

\startuseMPgraphic{page}
  StartPage ;
    def DrawShape (expr p, c) =
      fill p withcolor c ;
    enddef ;
    numeric ca, cb ; ca := 1 ; cb := .25 ;
    path pa, pb, pc, pd ;
    pa := boundingbox (center Page -- llcorner Page) ;
    pb := boundingbox (center Page -- lrcorner Page) ;
    pc := boundingbox (center Page -- urcorner Page) ;
    pd := boundingbox (center Page -- ulcorner Page) ;
    path pe, pf ;
    pe := .5[llcorner Page,lrcorner Page] --
          center pd randomized (.5*bbwidth(pd),.5*bbheight(pd)) --
          center pb randomized (.5*bbwidth(pb),.5*bbheight(pb)) --
          .5[ulcorner Page,urcorner Page] ;
    pf := .5[llcorner Page,ulcorner Page] --
          center pa randomized (.5*bbwidth(pa),.5*bbheight(pa)) --
          center pc randomized (.5*bbwidth(pc),.5*bbheight(pc)) --
          .5[lrcorner Page,urcorner Page] ;
    pa := llcorner Page -- pe -- ulcorner Page -- cycle ;
    pb := llcorner Page -- pf -- lrcorner Page -- cycle ;
    pc := urcorner Page -- reverse pe -- lrcorner Page -- cycle ;
    pd := ulcorner Page -- pf -- urcorner Page -- cycle ;
    DrawShape(pa,transparent(1,.5,\MPcolor{cyan})) ;
    DrawShape(pb,transparent(1,.5,\MPcolor{magenta})) ;
    DrawShape(pc,transparent(1,.5,\MPcolor{yellow})) ;
    DrawShape(pd,transparent(1,.5,\MPcolor{black})) ;
    interim linecap  := butt ;
    interim linejoin := mitered ;
    draw pe withpen pencircle scaled .5cm withcolor white ;
    draw pf withpen pencircle scaled .5cm withcolor white ;
    draw Page withpen pencircle scaled .5cm withcolor white ;
  StopPage ;
\stopuseMPgraphic

\setuphead
  [section]
  [style=\ss\bfb]

\startsetups [titlepage] 

  \definelayer
    [fullpage]
    [width=\paperwidth,
     height=\paperheight]
  
  \setlayer
    [fullpage]
    {\useMPgraphic{page}}
  
  \setlayerframed
    [fullpage]
    [corner={right,bottom},
     location={left,top},
     hoffset=2cm,
     voffset=1cm]
    [offset=overlay,
     frame=off,
     foregroundcolor=white]
    {\definedfont[SansBold at 2.5cm]%
     Color}
  
  \setlayerframed
    [fullpage]
    [corner={right,bottom},
     location={left,top},
     hoffset=1cm,
     voffset=1cm]
    [offset=overlay,
     frame=off,
     foregroundcolor=white]
    {\definedfont[SansBold at 1cm]%
     \rotate{\vbox{\hbox{Separation}\par\removedepth}}}
  
  \setlayerframed
    [fullpage]
    [corner={right,top},
     location={left,bottom},
     hoffset=1cm,
     voffset=1cm]
    [offset=overlay,
     frame=off,
     foregroundcolor=white]
    {\definedfont[SansBold at 1cm]%
     \rotate{\vbox{\hbox{Hans Hagen}\par\removedepth}}}
  
    \placelayer[fullpage] 
  
\stopsetups 

% \startuseMPgraphic{strip}{color}
%   linear_shade (unitsquare xyscaled(TextWidth/2,1cm),
%     5, \MPvariable{color}, white) ;
%   linear_shade (unitsquare xyscaled(TextWidth/2,1cm) 
%     shifted(TextWidth/2,0), 7, \MPvariable{color}, white) ;
% \stopuseMPgraphic

\startuseMPgraphic{strip}{color}
  linear_shade (unitsquare xyscaled(TextWidth,1cm), 5, 
    \MPvariable{color}, white) ;
\stopuseMPgraphic

\startuseMPgraphic{bar}{color}
  fill unitsquare xyscaled(TextWidth,1cm) withcolor \MPvariable{color} ;
\stopuseMPgraphic

\starttext 

\setuplayout[page] 

\startstandardmakeup \setups[titlepage] \stopstandardmakeup

\setuplayout[reset]

\section{Color spaces}

Currently \CONTEXT\ supports four color spaces: gray, \RGB,
\CMYK, and spot colors. The first three are defined rather
straightforward. The subtractive colors in \CMYK\ color
space are defined as follows:

\startbuffer
\definecolor[my-cyan]    [c=1]
\definecolor[my-magenta] [m=1]
\definecolor[my-yellow]  [y=1]
\definecolor[my-black]   [k=1]
\stopbuffer

\typebuffer \getbuffer

These colors (also known as process colors) show up as follows:

\startlinecorrection
  \useMPgraphic{strip}{color=my-cyan}
\stoplinecorrection

\startlinecorrection
  \useMPgraphic{strip}{color=my-magenta}
\stoplinecorrection

\startlinecorrection
  \useMPgraphic{strip}{color=my-yellow}
\stoplinecorrection

\startlinecorrection
  \useMPgraphic{strip}{color=my-black}
\stoplinecorrection

In the \RGB\ color space, we use additive colors:

\startbuffer
\definecolor[my-red]   [r=1]
\definecolor[my-green] [g=1]
\definecolor[my-blue]  [b=1]
\stopbuffer

\typebuffer \getbuffer

This time we get:

\startlinecorrection
\useMPgraphic{strip}{color=my-red}
\stoplinecorrection
\startlinecorrection
\useMPgraphic{strip}{color=my-green}
\stoplinecorrection
\startlinecorrection
\useMPgraphic{strip}{color=my-blue}
\stoplinecorrection

Gray colors are defined with:

\startbuffer
\definecolor[my-black] [s=0]
\stopbuffer

\typebuffer \getbuffer

Watch how this time we start at zero!

\startlinecorrection
\useMPgraphic{strip}{color=my-black}
\stoplinecorrection

When defining a color in \CMYK\ color space, you can mix
\type {c}, \type {m}, \type {y} and \type {k} components. 
On the cover you see four colors defines as:

\typebuffer[colors]

This leaves us spot colors. These are defined in two steps.
First you define the spot color itself:

\startbuffer
\definecolor[my-spot][c=.7,m=.5,y=.3,k=.1]
\stopbuffer

\typebuffer \getbuffer

As any other primary color, this color covers a range:

\startlinecorrection
\useMPgraphic{strip}{color=my-spot}
\stoplinecorrection

You can now define colors in this range:

\startbuffer
\definecolor[bright][my-spot][p=1]
\definecolor[dark]  [my-spot][p=.5]
\stopbuffer

\typebuffer \getbuffer

\startlinecorrection
\useMPgraphic{bar}{color=bright}
\stoplinecorrection
\startlinecorrection
\useMPgraphic{bar}{color=dark}
\stoplinecorrection

The chained definition is needed because devices, like
screen devices, that cannot handle spot colors and need to
have a fall back. 

\section{Special effects}

Those familiar with \METAFUN\ will know that it supports
shading. Actually the previous shades were produced with
code like:

\starttyping
\startMPcode
linear_shade (unitsquare xyscaled (TextWidth, 1cm), 5,
  \MPcolor {my-spot}, white) ;
\stopMPcode
\stoptyping

As you see here, we use \type {\MPcolor} to tell \METAPOST\
what color to use. Because \CONTEXT\ will handle the color
definitions when including the graphic in the document, the
proper transformations will take place and the required
(\PDF) resources will be embedded.

Both \CONTEXT\ and \METAFUN\ support transparent colors. As
you may expect, transparency works for all color spaces.
However, the place where you define them differs:

\startbuffer
\definecolor[redish] [r=.5,t=.5,a=1]
\definecolor[cyanish][c=.5,t=.5,a=1]
\definecolor[grayish][s=.5,t=.5,a=1]
\stopbuffer

\typebuffer \getbuffer

but:

\startbuffer
\definecolor[spottish][c=.3,m=.5,y=.7,k=.1]
\stopbuffer

\typebuffer \getbuffer

with:

\startbuffer
\definecolor[brightish][spottish][p=1,t=.5,a=1]
\definecolor[darkish]  [spottish][p=.5,t=.5,a=1]
\stopbuffer

\typebuffer \getbuffer

The following graphic demonstrates that these colors indeed
are transparent.

\startlinecorrection
\startMPcode
  path p ; p := unitsquare xyscaled (4.5cm,5mm) ;
  fill p shifted (0,0cm) withcolor \MPcolor{redish} ;
  fill p shifted (0,1cm) withcolor \MPcolor{cyanish} ;
  fill p shifted (0,2cm) withcolor \MPcolor{grayish} ;
  fill p shifted (0,3cm) withcolor \MPcolor{brightish} ;
  fill p shifted (0,4cm) withcolor \MPcolor{darkish} ;
  path p ; p := unitsquare xyscaled (5mm,4.5cm) ;
  fill p shifted (0cm,0) withcolor \MPcolor{redish} ;
  fill p shifted (1cm,0) withcolor \MPcolor{cyanish} ;
  fill p shifted (2cm,0) withcolor \MPcolor{grayish} ;
  fill p shifted (3cm,0) withcolor \MPcolor{brightish} ;
  fill p shifted (4cm,0) withcolor \MPcolor{darkish} ;
  currentpicture := currentpicture xysized (CurrentWidth,3cm) ;
\stopMPcode
\stoplinecorrection

\section{Options}

You can mix color spaces in one document. However, what
color space you use, depends on the medium on which the
document will be read. 

For screen documents, \RGB\ is often best, if only because
screen devices use these three color channels themselves.
All channels zero gives you black, and all channels one
produces white.

Printing houses often prefer \CMYK\ and|/|or spot colors.
They will split the document into the colors used for
printing, so in case of \CMYK\ only, there will be four
variants, one for each process color. If a document has only
one or two colors, spot colors will give good results, if
only because no screen has to be used. Also, the less
channels, the cheaper the print.

One danger luring with \CMYK\ is that some applications
produce output depending on the medium that will be used:
plain paper, glossy paper, colored paper etc. It's best to
let this kind of fine||tuning to the raster image processor
that will produce the plates, since that device knows best
what it is (in|)|capable off and what compensations are
needed.

You can control \CONTEXT's color mechanism in several ways
using the \type {\setupcolors} command. First of all, you
need to enable colors:

\starttyping
\setupcolors[state=start]
\stoptyping

By default all four methods are supported, but you can
change this behaviour:

\starttyping
\setupcolors[rgb=no,spot=no]
\stoptyping

This will force \CONTEXT\ to use \CMYK\ and gray scales
only (gray scale takes less code and will remap on the
black process color).

As long as you use \type {\MPcolor} in your \METAPOST\
graphics, you can be sure that those graphics will obey the
color settings. In the background, \CONTEXT\ will convert
colors defined in the disabled color spaces into those
enabled.

Although it can be tempting to use names for colors like
{\em red} or {\em blue}, it makes sense to use more
meaningful names, like {\em important} or {\em dangerous}.

\starttyping
\definecolor[dangerous][r=1]
\stoptyping

Even better is to define colors in an indirect way:

\starttyping
\definecolor[red][r=1]
\definecolor[dangerous][red]
\stoptyping

This means that when we map more colors on red, we only have
one definition of red to maintain. At the time of definition
of {\em dangerous}, only the name is mapped, which means
that we can change the definition of red afterwards. If you
want to freeze colors at definition time, say:

\starttyping
\setupcolors[expansion=yes]
\stoptyping

You can specify colors in your preferred color space, but 
your preferences may not match those of your printer or 
printing house. For that reason, \CONTEXT\ provides color 
conversion. You can for instance force \CMYK\ output by 
saying: 

\starttyping 
\setupcolors[rgb=no,spot=no,cmyk=yes]
\stoptyping 

If you also turn off \CMYK\ support, you will get weighted 
gray scales. You can also demand reduced \CMYK\ color, where
the black component is avoided, by setting \type {reduction}
to \type {yes}. \in {Table} [tab:values] shows what happens 
when such conversions are asked for. We use the previously 
defined colors.  

\startbuffer
\def\M#1{\setupcolors[rgb=yes,cmyk=yes,spot=yes]\showcolormodes\color[#1]{}}
\def\R#1{\setupcolors[rgb=yes,cmyk=no, spot=no]\showcolormodes\color[#1]{}}
\def\C#1{\setupcolors[rgb=no, cmyk=yes,spot=no]\showcolormodes\color[#1]{}}
\def\X#1{\setupcolors[rgb=no, cmyk=yes,spot=no,reduction=yes]\showcolormodes\color[#1]{}}
\def\G#1{\setupcolors[rgb=no, cmyk=no, spot=no]\showcolormodes\color[#1]{}}
\switchtobodyfont[small]
\starttabulate[|B||||||]
\HL
\NC color    \NC \bf mixed      \NC \bf rgb        \NC \bf cmy (reduced)
             \NC \bf cmyk       \NC \bf gray       \NC \NR
\HL
\NC red      \NC \M{my-red}     \NC \R{my-red}     \NC \X{my-red}
             \NC \C{my-red}     \NC \G{my-red}     \NC \NR
\NC green    \NC \M{my-green}   \NC \R{my-green}   \NC \X{my-green}
             \NC \C{my-green}   \NC \G{my-green}   \NC \NR
\NC blue     \NC \M{my-blue}    \NC \R{my-blue}    \NC \X{my-blue}
             \NC \C{my-blue}    \NC \G{my-blue}    \NC \NR
\NC cyan     \NC \M{my-cyan}    \NC \R{my-cyan}    \NC \X{my-cyan}
             \NC \C{my-cyan}    \NC \G{my-cyan}    \NC \NR
\NC magenta  \NC \M{my-magenta} \NC \R{my-magenta} \NC \X{my-magenta}
             \NC \C{my-magenta} \NC \G{my-magenta} \NC \NR
\NC yellow   \NC \M{my-yellow}  \NC \R{my-yellow}  \NC \X{my-yellow}
             \NC \C{my-yellow}  \NC \G{my-yellow}  \NC \NR
%\NC black   \NC \M{my-black}   \NC \R{my-black}   \NC \X{my-black}
%            \NC \C{my-black}   \NC \G{my-black}   \NC \NR
\NC spot     \NC \M{my-spot}    \NC \R{my-spot}    \NC \X{my-spot}
             \NC \C{my-spot}    \NC \G{my-spot}    \NC \NR
\NC bright   \NC \M{bright}     \NC \R{bright}     \NC \X{bright}
             \NC \C{bright}     \NC \G{bright}     \NC \NR
\NC dark     \NC \M{dark}       \NC \R{dark}       \NC \X{dark}
             \NC \C{dark}       \NC \G{dark}       \NC \NR
\HL
\stoptabulate
\stopbuffer

\placetable
  [here]
  [tab:values]
  {This table show what values end up in the output files 
   when conversion takes place.}
  {\getbuffer}

Due to the nature of \TEX\ and the control needed for
output devices, you need to be aware of interference of color
switches and typesetting. The low level color control
commands (for instance literal \PDF\ code) ends up in the
page stream. We will demonstrate this with an example (here 
we use the same cyan variant as on the cover): 

\startnarrower \showlowlevelstream \bgroup 
We start with no color, \color [cyan] {but now we have
switched to \currentcolor, \color [bright] {and now we're
typesetting in \currentcolor,} and some \currentcolor, and
\color [redish] {after a bit of \currentcolor}, we're back
at \currentcolor\ again.} 
\egroup \stopnarrower

The output (\DVI\ converted to \POSTSCRIPT, or \PDF) is
made up of so called independent pages. This means that,
apart from resources like fonts, printing a page should not
demand processing of previous pages. So, when we use
special features, like color, we have no knowledge about
the current color at a page break. Therefore, we need to
synchronize colors after a page break. 

\TEX\ invokes the output routine (page builder) after a
paragraph is processed. If the paragraph is typeset in a
color, the color is already reset to its orginal. However,
when \TEX\ decides that it should break in the middle of
that particular paragraph, we have to make sure that we add
codes like shown in the previous example at the top of the
next page. \CONTEXT\ keeps track of colors and tries to
synchronize the current state and the output state as good
as possible. 

Starting and ending a color obeys \TEX's grouping mechanism.
This means that when you want to typeset a document as a 
whole in a color, you need to group the whole document. This
can interfere with the final page break, and as a result, 
low level data streams like the ones shown before may end up
on a page of their own, which results in an empty page. In 
order to prevent such problems, you can set the document 
color with: 

\starttyping
\setupcolors[textcolor=somecolor]  
\stoptyping 

Similar complications can occur when you use colors in the
\type {\framed} command, which is also used extensively in
core macros. The \type {\color} command expects an argument,
and therefore exactly knows where to end coloring. Commands 
like \type {\red} on the other hand, hook themselves into 
the grouping mechanism. As a result, their behaviour may 
interfere with other mechanism, like automatic strut 
placement. Take the following text (here we use the same 
magenta variant as on the cover):

\startbuffer
The small bars shown here are called struts. Normally they
are invisible, but they play an important role in ensuring
consistent vertical spacing. You can {\magenta visualize} 
them with \type {\showstruts}.
\stopbuffer

\typebuffer

If we frame this text with:

\starttyping
\framed[align=middle]{\cyan...} 
\stoptyping

we get:

\startbaselinecorrection
\framed[align=middle]{\cyan\getbuffer}
\stopbaselinecorrection

At first sight the result looks ok, but let's visualize the 
struts:

\startbaselinecorrection \showstruts 
\framed[align=middle]{\cyan\getbuffer}
\stopbaselinecorrection

As you can see, there is some white space before the strut. 
Normally such a strut tries to get rid of it, so how come 
that it fails here? Let's also visualize the low level 
streams: 

\startbaselinecorrection \showstruts \showlowlevelstream 
\framed[align=middle]{\cyan\getbuffer}
\stopbaselinecorrection

It now becomes clear that the color commands end up before 
the strut, and as a result there is no way to go back. This 
is why in such situations it's better to apply the colors in
a different way: 

\starttyping 
\framed[align=middle,foregroundcolor=cyan]{...}
\stoptyping 

This results in:

\startbaselinecorrection \showstruts \showlowlevelstream 
\framed[align=middle,foregroundcolor=cyan]{\getbuffer}
\stopbaselinecorrection

Now, when we turn off the visualization, we get the properly
aligned text: 

\startbaselinecorrection 
\framed[align=middle,foregroundcolor=cyan]{\getbuffer}
\stopbaselinecorrection

The main message here is that, when you apply color, you 
need to be aware of the fact that this introduces (normally)
invisible content. Therefore, in case of doubt, use the 
\type {\color} or \type {\startcolor} command. 

\section{Groups and palets}

The reference manual describes in detail how you can group 
colors and how you can define color palets. Since we're 
dealing with color based workflows here, we only want to 
stress that using color palets may ease your task when one 
style has to serve multiple documents with different color 
settings. You define a palet as follows: 

\startbuffer
\definepalet[first] [maincolor=cyan,subcolor=yellow,extracolor=magenta]
\definepalet[second][maincolor=magenta,subcolor=cyan,extracolor=yellow]
\stopbuffer

\getbuffer \typebuffer

You can inherit them as well: 

\starttyping 
\definepalet [default] [first]
\stoptyping 

This means that you can isolate color definitions pretty 
well. Such a palet is enabled by: 

\starttyping 
\setuppalet[first]
\stoptyping

Because in the end, we still deal with colors, all the
color manipulations mentioned in this document apply to
palets as well. You can inspect a palet by \type
{\showpalet}, as demonstrated in \in {figure} [fig:palet].
The gray bars are the weighted gray alternatives. 

\placefigure
  [here] [fig:palet]
  {Two color palets visualized.}
  {\startcombination
     {\showpalet[first] [number]} {first} 
     {\showpalet[second][number]} {second}
   \stopcombination}

\section{Separation}

We now come to the main topic of this manual: color
separation. An \RGB\ color has three color components,
while an \CMYK\ color has four. Although modern printing
devices should be able to deal with both, in print the
\CMYK\ color space is preferred. 

When printing in full color, the four colors are printed 
separately, which comes down to printing four instances of 
the same document in succession. The separation needed for 
this can best be dealt with in the engine itself, but there 
may be occasions where you need to provide the separations 
yourself. 

You get the cyan channel version of a document when you 
say:

\starttyping 
\setupcolors[split=c]
\stoptyping 

In a similar way, you can ask for a magenta (\type {m}),
yellow \type {y}, red (\type {r}), green (\type {g}), blue
(\type{b)}, or black (\type {s}) version. If you want a spot
color, you set \type {split} to the desired color: 

\starttyping
\setupcolors[split=yourspot]
\stoptyping 

A more convenient way to do this, is to use \TEXEXEC: 

\starttyping
texexec --separation=c 
\stoptyping 

There is a preliminary script \type {texsplit} that
automatically splits the document, using information stored
in the auxiliary file \type {jobname.tuo}.  

Graphics are a complication in this process. For
black||and||white graphics, you can specify in which
channel they have to end up:  

\starttyping
\externalfigure[graphic][split=c]
\stoptyping 

Color graphics need to be separated independent of the 
document and we assume that this has been done. We discuss  
three ways to deal with including them. 

A convenient way to assign the (separated) graphics to a 
channel is the following:

\starttyping
\externalfigure[graphic\colorsplitsuffix]
\stoptyping 

In this case, the original graphic is called \type
{graphic}, while the cyan one is called \type {graphic-c},
and the rest accordingly. 

An alternative is to collect the graphics in a figure 
database and just switch databases. This saves you the 
effort of specifying the suffix each time. 

\starttyping
\usefigurebase[graphics\colorsplitsuffix]
\stoptyping

Instead of a suffix, you can apply a prefix (which is 
sometimes handier when you want to list files): 

\starttyping
\usefigurebase[\colorsplitprefix graphics]
\stoptyping

Alternatively, you can use a different path for the split 
graphics: 

\starttyping
\setupexternalfigures[directory={./figures\colorsplitsuffix}]
\stoptyping

A third alternative is to use modes. You can of course 
introduce special modes, like:

\starttyping
\startmode[cyan]
  \setupcolors[split=c]
  \usefigurebase[cyan]
\stopmode
\stoptyping

Here you need to enable a mode or run \TEXEXEC\ with \type
{--mode=cyan}. However, \CONTEXT\ provides channel specific
system modes. A system mode is prefixed by a \type {*}. 

\starttyping
\startmode[*color-c]
  all kind of channel specific settings 
\stopmode
\stoptyping 

All this may sound complicated, but color separation is a
complex matter anyway, so you'd better know what you're
doing. This is especially true when you set \type
{criterium}, which by default is set to \type {all}. This
means that when a 100\% spot color is used, it is assumed
to completely knock out its background. On the other hand,
for instance a 50\% spot color is painted on a white
background. 

\starttyping
\setupcolors[criterium=none] 
\stoptyping 

When \type {criterium} is set this way, the background is
always painted in white. This assumes a rather precise
printing workflow. 

For external figures, you can specify a background color 
with \type {splitcolor}. By default this color is set to 
white. 

\starttyping 
\setupexternalfigures[splitcolor=white]
\stoptyping 

\section{Bonus example}

When color separation was added \CONTEXT, some experiments 
were conducted with coloring bitmap graphics. One of the 
poor||mans colorizations is presented here. 

\startbuffer
\definecolor [p-one]   [c=1,y=.3,k=.3]
\definecolor [p-two]   [c=.25,y=.8]
\definecolor [p-three] [c=.9,y=.15]

\definecolor [one-1]   [p-one]   [p=1,t=1,a=overlay]
\definecolor [one-2]   [p-one]   [p=1,t=1,a=softlight]

\definecolor [two-1]   [p-two]   [p=1,t=1,a=overlay]
\definecolor [two-2]   [p-two]   [p=1,t=1,a=softlight]

\definecolor [three-1] [p-three] [p=1,t=1,a=overlay]
\definecolor [three-2] [p-three] [p=1,t=1,a=softlight]
\stopbuffer

First we define a few transparent colors. Out of the 
alternatives, we choose \type {overlay} and \type 
{softlight}. (More about transparency can be found in the 
\METAFUN\ manual.) 

\getbuffer \typebuffer

\startbuffer[mill]
\startcombination[6*1]
  {\externalfigure[mill.png][width=.125\textwidth,
     split=p-two,foregroundcolor=two-1]}     {\white overlay}
  {\externalfigure[mill.png][width=.125\textwidth,
     split=p-two,foregroundcolor=two-2]}     {\white softlight}
  {\externalfigure[mill.png][width=.125\textwidth,
     split=p-one,foregroundcolor=one-1]}     {\white overlay}
  {\externalfigure[mill.png][width=.125\textwidth,
     split=p-one,foregroundcolor=one-2]}     {\white softlight}
  {\externalfigure[mill.png][width=.125\textwidth,
     split=p-three,foregroundcolor=three-1]} {\white overlay}
  {\externalfigure[mill.png][width=.125\textwidth,
     split=p-three,foregroundcolor=three-2]} {\white softlight}
\stopcombination
\stopbuffer

\typebuffer

\startbuffer[back]
\definecolor[darkgray][s=.3]

\defineframedcontent
  [mill]
  [background=color,
   backgroundcolor=darkgray,
   frame=off,
   offset=1ex]
\stopbuffer

\getbuffer[back]

This cheap trick leads to: 

\startbaselinecorrection
\startframedcontent[mill] \getbuffer[mill] \stopframedcontent
\stopbaselinecorrection

In this example, the gray background was accomplished with: 

\typebuffer[back]

and: 

\starttyping 
\startframedcontent[mill] ... \stopframedcontent
\stoptyping 

This example was part of a test document with three
spot colors and black, which demands four splits. The
original and its four splits are shown on the next pages. 

\placefigure
  {A test document with three spot colors \& black.} 
  {\externalfigure[septest.pdf][frame=on,page=1,height=.9\textheight]}

\page \combinepages[septest.pdf][nx=2,ny=2,start=2,stop=5,frame=on]

\stoptext
